% ============================================================================
% Fichier  : partie1/P01_C05_etat_art.tex
% Titre    : Etat de l'Art et Defis Contemporains
% Auteur   : MINKA MI NGUIDJOÏ Thierry Emmanuel
% Version  : 1.0 -- Juin 2026
% Statut   : REDIGE
% Sources  : 4_etat_art.tex, 20_benchmarking.tex, M6_OSINT.docx,
%            6Competences_Techniques_1.PDF, minka2025cro
% ============================================================================

\chapter{État de l'Art et Défis Contemporains}
\label{chap:etat-art}

\epigraph{%
    « La science progresse en faisant danser les faits aux rythmes
    de nouvelles théories. »%
}{--- Marcel Proust}

\niveauChapitre{Fondamental}{%
    Ce chapitre est le dernier de la Partie~I. Il mobilise les concepts
    des chapitres précédents pour cartographier l'état de la discipline
    en 2025--2026 et projeter son évolution à l'horizon 2030.%
}

\begin{boxobjectifs}
\begin{itemize}
    \item Situer les cinq frontières actives de l'investigation numérique
          en 2025--2026 et comprendre les défis spécifiques de chacune.
    \item Analyser la valeur et les limites de l'IA forensique~: LLM pour
          triage, machine learning pour classification de malware,
          deepfake detection.
    \item Comprendre les défis forensiques propres au cloud, à l'IoT et
          au mobile dans le contexte africain et camerounais.
    \item Maîtriser les fondamentaux de l'OSINT légal pour les OPJ~:
          outils, cadre juridique, limites.
    \item Utiliser le framework DICES pour positionner les pratiques
          camerounaises par rapport aux standards mondiaux.
    \item Construire une roadmap de développement professionnel réaliste
          sur 36~mois dans le contexte camerounais.
\end{itemize}
\end{boxobjectifs}

\bigskip

% ============================================================================
\section*{Bilan d'une Discipline en Accélération}
% ============================================================================

La Partie~I de ce cours a posé les fondements~: philosophie, histoire,
affaires fondatrices, théories mathématiques. Ce dernier chapitre de la
Partie~I change de registre~: il décrit ce qui se passe \emph{maintenant},
les défis que la discipline affronte \emph{aujourd'hui}, et les compétences
que vous devrez avoir construites \emph{dans cinq ans}.

\medskip

L'investigation numérique de 2025 n'est plus celle de 2015. Cinq ruptures
simultanées la reconfigurent en profondeur~:
\begin{enumerate}
    \item L'\textbf{IA} entre dans la chaîne d'analyse --- comme outil
          d'accélération du triage et comme vecteur d'attaque sophistiqué.
    \item Les \textbf{deepfakes} brouillent la frontière entre le vrai
          et le fabriqué, forçant la discipline à développer une
          sous-spécialité d'authentification des médias.
    \item Le \textbf{cloud et l'IoT} pulvérisent la notion de scène de
          crime localisée~: les preuves sont partout, souvent éphémères,
          toujours multi-juridictionnelles.
    \item L'\textbf{OSINT} s'est professionnalisé au point d'être une
          compétence à part entière, avec ses outils, son droit et son
          éthique propres.
    \item La \textbf{cryptographie post-quantique} est en cours de
          déploiement~: l'investigateur de demain doit comprendre ce que
          cela change pour la chaîne de preuves.
\end{enumerate}

\separateur

% ============================================================================
\section{IA Forensique~: Puissance et Limites}
\label{sec:ia-forensique}
% ============================================================================

\subsection{L'IA comme accélérateur de triage}

Le triage forensique --- identifier en quelques minutes les artefacts
pertinents parmi des milliers --- est la tâche où l'IA apporte le gain
le plus immédiat. Les \termecle{LLM}\index{LLM forensique}
(\textit{Large Language Models}) peuvent~:
\begin{itemize}
    \item \textbf{Résumer et catégoriser} des volumes d'emails, de journaux
          ou de logs qui dépassent la capacité d'analyse humaine en temps
          raisonnable.
    \item \textbf{Identifier des anomalies linguistiques} dans des
          communications --- syntaxe inhabituelle, traduction automatique,
          incohérences de registre.
    \item \textbf{Générer des hypothèses} à partir d'un corpus de preuves
          hétérogènes~: un LLM peut proposer des connections entre artefacts
          que l'analyste humain n'aurait pas immédiatement vus.
\end{itemize}

\subsection{Classification automatique de malware}

Les approches de machine learning pour la classification de malware
reposent sur trois familles de caractéristiques~:

\begin{tableaucours}{p{3.5cm}p{3.5cm}p{4.5cm}}{%
    \tableauHeader{Approche} &
    \tableauHeader{Caractéristiques} &
    \tableauHeader{Forces / Limites}
}
Statique & Chaînes, imports, sections PE, entropie &
    Rapide, pas d'exécution~; contourné par l'obfuscation \\
Dynamique & Appels système, trafic réseau, comportement &
    Robuste à l'obfuscation~; lent, risque d'environnement sandbox \\
Hybride & Combinaison statique + dynamique &
    Meilleure précision~; coûteux en calcul \\
\end{tableaucours}
\vspace{-0.8em}
\begin{center}{\small\itshape Tableau~5.1 --- Approches ML pour la classification de malware}
\end{center}

\subsection{Les limites de l'IA forensique --- la crise d'opposabilité}

L'IA forensique pose une question CRO fondamentale~: si une conclusion
repose sur un modèle opaque, comment la défendre devant un tribunal~?

\begin{boxCRO}
\textbf{Analyse CRO --- IA forensique}

$\mathcal{R}$ (Fiabilité)~: les modèles de ML bien entraînés atteignent
des précisions de 95--99\% sur des jeux de données équilibrés.

$\mathcal{O}$ (Opposabilité)~: \textbf{c'est la crise.} Un avocat de la
défense peut demander l'explication complète du modèle. Si la réponse
est ``le réseau de neurones a attribué une probabilité de 0.97''
sans explication compréhensible, le juge peut refuser la conclusion.
L'arrêt \textit{Loomis v. Wisconsin} (2016) aux États-Unis illustre
ce risque pour les algorithmes de prédiction en justice.

$\mathcal{C}$ (Confidentialité)~: les modèles entraînés sur des données
réelles (cas judiciaires précédents) peuvent mémoriser des informations
sur des affaires passées et les révéler via des attaques par inversion.

\cro{$\downarrow$ risque de fuite}{$\uparrow$ précision technique}{$\downarrow$ crise d'explicabilité}

\textbf{Règle pratique~:} utiliser l'IA pour le \emph{triage} (hypothèses,
orientation), jamais comme \emph{preuve primaire}. L'analyste humain
confirme et documente chaque conclusion de l'IA.
\end{boxCRO}

\begin{boxalert}
\textbf{Règle d'or~: l'IA oriente, l'humain constate.}

En contexte judiciaire camerounais, aucun texte n'autorise encore une
preuve basée uniquement sur un algorithme opaque. L'investigateur doit
systématiquement~: (1)~documenter l'outil utilisé (nom, version,
paramètres), (2)~vérifier manuellement les conclusions de l'IA sur un
échantillon, (3)~présenter la conclusion IA comme \emph{indication}
corroborée par une analyse humaine, jamais comme preuve autonome.
\end{boxalert}

% ============================================================================
\section{Deepfake Detection~: Authentifier les Médias Numériques}
\label{sec:deepfakes}
% ============================================================================

\subsection{Le défi de l'authenticité des médias}

Les \termecle{deepfakes}\index{deepfake} --- vidéos, images, audios générés
ou manipulés par IA --- constituent l'une des menaces forensiques les plus
graves de la décennie. Trois catégories~:
\begin{itemize}
    \item \textbf{Face swap}~: substitution du visage d'une personne par
          celui d'une autre dans une vidéo existante.
    \item \textbf{Voice cloning}~: reproduction convaincante de la voix
          d'une personne à partir d'un échantillon de quelques secondes.
    \item \textbf{Document forgery}~: génération de documents officiels
          (relevés bancaires, identités, contrats) cohérents avec les
          templates réels.
\end{itemize}

\subsection{Méthodes de détection}

Quatre approches complémentaires~:
\begin{enumerate}
    \item \textbf{Analyse de cohérence physique}~: les deepfakes génèrent
          souvent des incohérences dans les reflets des yeux, le
          mouvement des oreilles, la synchronisation labiale.
    \item \textbf{Analyse spectrale}~: les artefacts de compression GAN
          (\textit{Generative Adversarial Network}) laissent des
          signatures dans le domaine fréquentiel (FFT) détectables
          même après recompression.
    \item \textbf{Métadonnées et provenance}~: la
          \termecle{C2PA}\index{C2PA} (\textit{Coalition for Content
          Provenance and Authenticity}) développe un standard de
          signature cryptographique à la capture~: chaque photo prise
          par un appareil certifié porte une signature vérifiable.
    \item \textbf{Analyse comportementale}~: les modèles de langage
          génèrent des textes avec des patterns statistiques spécifiques
          détectables par des outils comme GPTZero ou Binoculars.
\end{enumerate}

\begin{boxartefact}{Vérification d'authenticité d'image -- commandes pratiques}
\textbf{Étape 1 --- Extraction des métadonnées EXIF et C2PA~:}
\begin{lstlisting}[language=bash_forensique]
# ExifTool -- extraction complete
exiftool -a -u -g1 image_suspecte.jpg
# Chercher : Make, Model, GPS, CreateDate, Software
# ALERTE : 'Software: GIMP' sur une photo prétendument originale
# ALERTE : absence de champ GPS sur une photo censée être localisée

# Analyse des erreurs de niveau JPEG (ELA -- Error Level Analysis)
# Disponible via fotoforensics.com ou imagemagick
convert image.jpg -auto-level ela_result.jpg
# Les zones retouchées apparaissent avec un niveau d'erreur différent
\end{lstlisting}

\textbf{Étape 2 --- Analyse spectrale (FFT)~:}
\begin{lstlisting}[language=bash_forensique]
python3 -c \"
import cv2, numpy as np, matplotlib.pyplot as plt
img = cv2.imread('suspect.jpg', 0)
f = np.fft.fft2(img)
fshift = np.fft.fftshift(f)
magnitude = 20*np.log(np.abs(fshift)+1)
plt.imshow(magnitude, cmap='gray')
plt.savefig('spectrum.png')
\"
# Un spectre en croix ou en échiquier indique une génération GAN
\end{lstlisting}
\end{boxartefact}

\begin{boxjuris}
\textbf{Statut juridique des deepfakes au Cameroun~:}

La \loicam{2010/012}, art.~72 (fraude informatique) et art.~71
(modification de données) couvrent la création de documents falsifiés,
y compris les deepfakes à des fins frauduleuses. La \loicam{2024/017},
art.~59 (utilisation frauduleuse des données personnelles) s'applique
lorsque le deepfake utilise le visage ou la voix d'une personne sans
son consentement.
\end{boxjuris}

% ============================================================================
\section{Forensique Cloud, IoT et Mobile~: la Scène de Crime Distribuée}
\label{sec:cloud-iot-mobile}
% ============================================================================

\subsection{Les trois défis structurels du cloud forensique}

La forensique cloud n'est pas une extension de la forensique système~:
c'est une discipline qui doit être repensée. Trois défis fondamentaux~:

\begin{enumerate}
    \item \textbf{Multitenancy}~: les données d'un client sont physiquement
          colocalisées avec celles d'autres clients sur les mêmes serveurs.
          L'investigateur ne peut pas simplement ``saisir le serveur''.

    \item \textbf{Éphéméralité}~: les instances cloud (containers,
          fonctions serverless) durent quelques secondes à quelques minutes.
          Les logs sont supprimés selon des politiques de rotation que
          l'investigateur ne contrôle pas.

    \item \textbf{Multi-juridiction}~: un service hébergé chez AWS peut
          avoir ses données réparties entre des datacenters en Irlande,
          aux États-Unis et à Singapour. Quelle loi s'applique~?
          Le \textit{CLOUD Act} américain (2018) donne aux autorités US
          accès aux données hébergées par des entreprises américaines
          \emph{n'importe où dans le monde}~; cela entre en conflit avec
          la souveraineté de nombreux États.
\end{enumerate}

\begin{boxoutils}{Forensique cloud --- commandes de collecte AWS/Azure}
\textbf{AWS CloudTrail --- journaux d'audit~:}
\begin{lstlisting}[language=bash_forensique]
# Lister les événements récents CloudTrail (30 derniers jours)
aws cloudtrail lookup-events \\
    --lookup-attributes AttributeKey=EventName,AttributeValue=DeleteBucket \\
    --max-results 50

# Exporter les logs d'un bucket S3 spécifique vers l'investigation
aws s3 cp s3://bucket-suspect/logs/ ./evidence/logs/ --recursive
# Calculer les hash immédiatement après téléchargement
sha256sum ./evidence/logs/*
\end{lstlisting}

\textbf{Azure Activity Log~:}
\begin{lstlisting}[language=bash_forensique]
# Lister les activités suspectes sur un abonnement Azure
az monitor activity-log list \\
    --start-time 2026-03-01 \\
    --end-time 2026-03-15 \\
    --query '[?level==`Critical`]' \\
    --output table
\end{lstlisting}
\end{boxoutils}

\subsection{IoT forensique~: l'explosion des surfaces de collecte}

Selon Cisco, plus de 29~milliards d'appareils IoT seront connectés en 2030.
Chacun est potentiellement une scène de crime~: routeur compromis, caméra IP,
montre connectée, véhicule, thermostat intelligent. Les défis~:

\begin{itemize}
    \item \textbf{Hétérogénéité}~: chaque fabricant utilise un système
          d'exploitation propriétaire. Les outils forensiques génériques
          (Autopsy, FTK) ne fonctionnent pas sur un firmware TP-Link.
    \item \textbf{Volatilité extrême}~: beaucoup d'appareils IoT ne
          conservent leurs logs qu'en RAM ou dans un buffer circulaire
          de quelques heures.
    \item \textbf{Extraction physique}~: certains appareils exigent
          une extraction via UART, JTAG ou chip-off --- des techniques
          matérielles qui nécessitent un équipement spécialisé.
\end{itemize}

\begin{boxscenario}{Caméra IP dans une affaire d'homicide --- Yaoundé, fictif}
Une caméra de vidéosurveillance IP a enregistré l'entrée d'un suspect dans
un immeuble le soir d'un homicide. L'immeuble utilise un système Hikvision.
La caméra est connectée à un NVR (enregistreur réseau).

\textbf{Procédure forensique~:}
\begin{enumerate}
    \item Identifier le modèle du NVR et son interface d'administration.
    \item Extraire les séquences vidéo via l'interface web ou via
          l'outil dédié (IVMS-4200 pour Hikvision).
    \item Calculer le hash de chaque fichier vidéo extrait.
    \item Documenter les paramètres de la caméra (heure système,
          synchronisation NTP~; un décalage horaire fausse tout
          témoignage basé sur la vidéo).
\end{enumerate}

\textbf{Attention CRO~:} \cro{$\sim$}{$\uparrow$ si hash documenté}{$\sim$ dépend horloge}
Le décalage horaire est le piège le plus fréquent des vidéos de
surveillance~: sans synchronisation NTP documentée, l'heure affichée
peut être fausse de plusieurs heures.
\end{boxscenario}

\subsection{Mobile forensique~: spécificités africaines}

En Afrique subsaharienne, le mobile est le premier vecteur de services
numériques (banque, communication, commerce). La forensique mobile y est
donc disproportionnellement centrale. Quatre niveaux d'acquisition selon
\norme{NIST SP~800-101r1}~:

\begin{tableaucours}{p{3.0cm}p{3.0cm}p{5.5cm}}{%
    \tableauHeader{Niveau} &
    \tableauHeader{Technique} &
    \tableauHeader{Données accessibles}
}
Niveau~1 & Extraction manuelle & Photos, contacts, SMS visibles à l'écran \\
Niveau~2 & Extraction logique (ADB) &
    Fichiers accessibles du système de fichiers, bases SQLite applications \\
Niveau~3 & Extraction physique &
    Image complète de la mémoire flash~; données supprimées récupérables \\
Niveau~4 & Chip-off / JTAG &
    Extraction directe de la puce mémoire~; dernier recours \\
\end{tableaucours}
\vspace{-0.8em}
\begin{center}{\small\itshape Tableau~5.2 --- Niveaux d'acquisition mobile}
(NIST SP~800-101r1)}\end{center}

% ============================================================================
\section{OSINT~: Renseignement en Source Ouverte}
\label{sec:osint}
% ============================================================================

\subsection{Définition, périmètre et cadre légal camerounais}

L'\termecle{OSINT}\index{OSINT} (\textit{Open Source Intelligence}) désigne
la collecte et l'analyse d'informations issues de sources légalement
accessibles au public~: sites web, réseaux sociaux, bases WHOIS, registres
officiels, archives, publications scientifiques.

\begin{boxjuris}
\textbf{Cadre légal OSINT au Cameroun~:}
\begin{itemize}
    \item Art.~41 \loicam{2010/012}~: droit au respect de la vie privée.
          L'OSINT ne couvre que les informations \emph{volontairement}
          rendues publiques par leur sujet.
    \item Art.~9 \loicam{2024/017}~: la collecte systématique (fichage)
          et le croisement avec des données non publiques peuvent
          entrer en conflit avec le principe de finalité déterminée.
    \item \textbf{Règle pratique~:} documenter dans chaque rapport OSINT
          la source exacte et la date de consultation. Une capture d'écran
          horodatée est indispensable --- un lien public aujourd'hui peut
          être privé à l'audience.
    \item \textbf{Infiltration en ligne~:} se faire passer pour quelqu'un
          d'autre en ligne pour obtenir des informations constitue une
          provocation ou surveillance couverte~; elle requiert une
          autorisation judiciaire et sort du cadre OSINT.
\end{itemize}
\end{boxjuris}

\subsection{Le cycle OSINT pour les OPJ}

Le cycle OSINT structuré en cinq étapes~:
\begin{enumerate}
    \item \textbf{Planification}~: définir la cible et les questions
          clés. Qu'est-ce que j'ignore et que je dois trouver~?
    \item \textbf{Collecte}~: identifier et exploiter les sources
          pertinentes (réseaux sociaux, WHOIS, enregistrements publics).
    \item \textbf{Traitement}~: organiser, filtrer, supprimer les doublons.
    \item \textbf{Analyse}~: recouper, vérifier, tirer des conclusions.
    \item \textbf{Diffusion}~: produire un rapport utilisable par le
          magistrat instructeur.
\end{enumerate}

\subsection{Boîte à outils OSINT pour les OPJ camerounais}

\begin{tableaucours}{p{2.8cm}p{2.5cm}p{6.0cm}}{%
    \tableauHeader{Outil} &
    \tableauHeader{Usage principal} &
    \tableauHeader{Application forensique}
}
Sherlock & Recherche de pseudo & Identifier tous les comptes d'une personne
    à partir d'un username (cas WOURI~: \texttt{duke\_cyber\_ng}) \\
WHOIS / dig & Infrastructure & Propriétaire d'un domaine~; serveurs DNS~;
    enregistrement d'un domaine de phishing \\
crt.sh & Certificats SSL & Trouver tous les sous-domaines d'un domaine~;
    identifier des sites de phishing (ex~: \texttt{*.uba-cameroun.com}) \\
dnstwist & Typosquatting & Générer et vérifier les variantes d'un domaine
    légitime \\
Shodan & Appareils exposés & Identifier des caméras, serveurs, IoT
    vulnérables associés à une cible \\
ExifTool & Métadonnées & GPS dans les photos, auteur, appareil~;
    localiser un suspect via ses publications Instagram \\
Blockchair & Blockchain & Tracer les transactions Bitcoin~;
    identifier les adresses de mixeurs (cas WOURI) \\
Maltego & Graphe de relations & Cartographier visuellement les connexions
    entre entités (personnes, domaines, IPs) \\
\end{tableaucours}
\vspace{-0.8em}
\begin{center}{\small\itshape Tableau~5.3 --- Boîte à outils OSINT pour les OPJ}
\end{center}

\begin{lstlisting}[language=bash_forensique,
    caption={Workflow OSINT complet --- cible fictive KOUAME Narcisse}]
# 1. Recherche de pseudo identifié (@koffi_biz_225)
sherlock koffi_biz_225
# -> identifie comptes Telegram, Instagram, Twitter/X

# 2. Reverse image search via exiftool si photo disponible
exiftool photo_profil.jpg | grep -i 'GPS\|Make\|Date'

# 3. Infrastructure : domaine utilisé pour le phishing
whois medequip-portugal.com
dig medequip-portugal.com ANY
crt.sh : medequip-portugal.com  # via navigateur

# 4. Adresses Bitcoin du réseau
# -> Blockchair.com : entrer l'adresse, visualiser le graphe de transactions
# -> Identifier les exchanges (Binance, Kraken) : portail LEA disponible

# 5. Maltego : cartographier les connexions
# Entités : KOUAME, @koffi_biz_225, +225 07XXXXXX,
#           medequip-portugal.com, 3 adresses BTC
\end{lstlisting}

% ============================================================================
\section{Benchmarking Mondial~: le Framework DICES}
\label{sec:dices}
% ============================================================================

\subsection{Le framework DICES}

Le framework \termecle{DICES}\index{DICES, framework} (\textbf{D}octrine,
\textbf{I}nfrastructure, \textbf{C}apacités, \textbf{E}cosystème,
\textbf{S}tratégie) permet de comparer les pratiques forensiques de
différentes juridictions sur cinq axes \cite{dfrws2001}~:
\begin{itemize}
    \item \textbf{D}octrine~: philosophie et approche conceptuelle
          (empirisme pragmatique ou rigueur formelle~?)
    \item \textbf{I}nfrastructure~: moyens techniques et organisationnels
          (laboratoires, équipements, plateformes cloud)
    \item \textbf{C}apacités~: compétences humaines et processus
          (formations, certifications, procédures standardisées)
    \item \textbf{E}cosystème~: environnement juridique et institutionnel
          (lois, tribunaux, organismes de régulation)
    \item \textbf{S}tratégie~: vision prospective et adaptation
          (investissements R\&D, coopération internationale)
\end{itemize}

\subsection{Positionnement comparatif}

\begin{tableaucours}{p{2.4cm}p{2.0cm}p{2.0cm}p{2.0cm}p{2.0cm}p{2.0cm}}{%
    \tableauHeader{Axe DICES} &
    \tableauHeader{USA/FBI} &
    \tableauHeader{Europe} &
    \tableauHeader{Afrique du Sud} &
    \tableauHeader{Nigeria} &
    \tableauHeader{Cameroun}
}
Doctrine & Empirisme & Droit civil & Mixte & Common law & Droit civil fr. \\
Infrastructure & $\bullet\bullet\bullet\bullet\bullet$ & $\bullet\bullet\bullet\bullet$ &
    $\bullet\bullet\bullet$ & $\bullet\bullet$ & $\bullet\bullet$ \\
Capacités & $\bullet\bullet\bullet\bullet\bullet$ & $\bullet\bullet\bullet\bullet$ &
    $\bullet\bullet\bullet$ & $\bullet\bullet$ & $\bullet\bullet$ \\
Ecosystème & CFAA + FRE & RGPD + Budapest & RICA & Budapest & Loi 2010 \\
Stratégie & IA + quantique & ENISA + C2PA & SAPS-DCI & EFCC & ANTIC \\
\end{tableaucours}
\vspace{-0.8em}
\begin{center}{\small\itshape Tableau~5.4 --- Comparaison DICES~: $\bullet$ = niveau de maturité}
\end{center}

\subsection{Atouts et lacunes du Cameroun~: une lecture honnête}

\textbf{Atouts réels~:}
\begin{itemize}
    \item Un cadre législatif complet et récent~: \loicam{2010/012} et
          \loicam{2024/017} couvrent les infractions numériques et la
          protection des données.
    \item L'ANTIC, institution fonctionnelle et présente dans les enquêtes.
    \item Une communauté technique en formation rapide (ENSP, Universités,
          entreprises privées).
    \item Une expérience terrain sur les fraudes MoMo~: le Cameroun forme
          ses propres enquêteurs sur des cas réels, pas des simulations.
\end{itemize}

\textbf{Lacunes à combler~:}
\begin{itemize}
    \item Absence de laboratoire forensique accrédité au niveau national.
    \item Faible infrastructure de partage de renseignements cyber
          (\textit{threat intelligence}) entre les institutions.
    \item Jurisprudence numérique naissante~: peu de décisions de justice
          commentées et publiées sur la preuve numérique.
    \item Ressources d'investigation en langue française et en contexte
          africain~: c'est précisément le vide que ce cours est conçu pour
          combler.
\end{itemize}

% ============================================================================
\section{Roadmap de Développement Professionnel}
\label{sec:roadmap}
% ============================================================================

\subsection{Les certifications de référence mondiales}

\begin{tableaucours}{p{3.0cm}p{2.5cm}p{3.0cm}p{3.0cm}}{%
    \tableauHeader{Certification} &
    \tableauHeader{Niveau} &
    \tableauHeader{Spécialité} &
    \tableauHeader{Coût estimé}
}
GIAC GCFE & Intermédiaire & Forensique Windows & 7~000~USD \\
GIAC GCFA & Avancé & Incident Response & 7~000~USD \\
GIAC GREM & Expert & Malware RE & 7~000~USD \\
EnCE & Intermédiaire & EnCase (outil) & 3~000~USD \\
OSCP & Avancé & Pentest & 1~500~USD \\
CompTIA CySA+ & Intermédiaire & Analyse cyber & 400~USD \\
Cellebrite CCPA & Intermédiaire & Mobile forensique & Variable \\
\end{tableaucours}
\vspace{-0.8em}
\begin{center}{\small\itshape Tableau~5.5 --- Certifications forensiques de référence
(d'après \cite{competences2024})}\end{center}

\subsection{Roadmap réaliste sur 36 mois en contexte camerounais}

\begin{tableaucours}{p{2.5cm}p{2.2cm}p{7.0cm}}{%
    \tableauHeader{Phase} &
    \tableauHeader{Durée} &
    \tableauHeader{Objectifs et activités}
}
Phase~1 --- Socle & Mois 1--6 &
    Maîtrise de Kali Linux, Autopsy, FTK~Imager, Volatility~3.
    Réalisation de 20~scénarios de laboratoire (images forensiques fournies).
    Lecture de ce cours. Obtention CompTIA CySA+ si financement disponible. \\
Phase~2 --- Spécialisation & Mois 7--18 &
    Choisir une spécialité~: (a)~forensique judiciaire (droit +
    procédure), (b)~réponse aux incidents (cloud + SIEM),
    (c)~forensique mobile (Android/iOS). Participation aux CTF forensiques
    (CyberDefenders, BlueTeamLabs). Contribution à 3~investigations réelles
    (stage ou DGSN). \\
Phase~3 --- Excellence & Mois 19--36 &
    Préparation GIAC GCFE ou GCFA. Rédaction d'un article technique
    sur un cas camerounais. Contribution à l'élaboration de procédures
    standardisées pour l'ANTIC ou la DGSN. \\
\end{tableaucours}
\vspace{-0.8em}
\begin{center}{\small\itshape Tableau~5.6 --- Roadmap 36~mois en contexte camerounais}
\end{center}

\begin{boxPNC}
\textbf{Ressources gratuites accessibles depuis le Cameroun~:}
\begin{itemize}
    \item \textbf{CyberDefenders.org}~: laboratoires forensiques avec
          images disque, captures réseau, fichiers mémoire. Niveaux
          Débutant à Expert. Gratuit.
    \item \textbf{BlueTeamLabsOnline.com}~: scénarios de réponse aux
          incidents. Gratuit pour les niveaux de base.
    \item \textbf{HackTheBox / DFIR Track}~: challenge forensique mensuel.
    \item \textbf{SANS Cyber Aces}~: cours de fondamentaux gratuit.
    \item \textbf{GitHub/volatility3}~: documentation officielle avec
          exemples d'analyse mémoire.
    \item \textbf{Journal of Digital Forensics, Security and Law}~:
          accès ouvert~; articles académiques en forensique.
\end{itemize}
\end{boxPNC}

% ============================================================================
\section{Synthèse~: la Carte de la Discipline en 2025--2026}
\label{sec:synthese-etat-art}
% ============================================================================

Ce chapitre clôt la Partie~I. La figure ci-dessous synthétise la carte
de la discipline telle qu'elle se présente en 2025--2026~:

\begin{figure}[H]
\centering
\begin{tikzpicture}[
    domaine/.style={
        draw=CouleurPrimaire, fill=CouleurDef, rounded corners=5pt,
        minimum width=3.0cm, minimum height=0.85cm,
        font=\small\bfseries\color{CouleurPrimaire}, align=center
    },
    defi/.style={
        draw=CouleurBordAlerte, fill=CouleurAlerte, rounded corners=4pt,
        minimum width=2.8cm, minimum height=0.75cm,
        font=\footnotesize\color{CouleurBordAlerte}, align=center
    },
    coeur/.style={
        draw=CouleurAccent, fill=CouleurAccent!20, rounded corners=6pt,
        minimum width=3.5cm, minimum height=1.1cm,
        font=\small\bfseries\color{CouleurAccent!80!black}, align=center
    },
    fl/.style={-{Stealth[length=4pt]}, thick, color=CouleurSecondaire}
]
% Centre : CRO
\node[coeur] (cro) at (0,0) {Trilemme CRO\\$\mathcal{C}$--$\mathcal{R}$--$\mathcal{O}$};
% Domaines actifs
\node[domaine] (ia)    at (-4,  2.5) {IA Forensique};
\node[domaine] (df)    at ( 0,  3.5) {Deepfake\\Detection};
\node[domaine] (cloud) at ( 4,  2.5) {Cloud / IoT\\Mobile};
\node[domaine] (osint) at ( 4, -2.5) {OSINT};
\node[domaine] (pq)    at (-4, -2.5) {Post-Quantique};
% Défis
\node[defi] (opp)  at (-2.5,  0.8) {Opposabilité IA~?};
\node[defi] (jur)  at ( 2.5,  0.8) {Juridiction\\cloud~?};
\node[defi] (hndl) at (-2.5, -0.8) {HNDL};
\node[defi] (priv) at ( 2.5, -0.8) {Vie privée\\vs enquête};
% Flèches
\draw[fl] (ia)    -- (cro);
\draw[fl] (df)    -- (cro);
\draw[fl] (cloud) -- (cro);
\draw[fl] (osint) -- (cro);
\draw[fl] (pq)    -- (cro);
\end{tikzpicture}
\caption{Carte de la forensique numérique en 2025--2026~: cinq frontières
actives toutes soumises aux tensions du Trilemme CRO}
\end{figure}

\separateur

\begin{boxresume}
\textbf{Ce qu'il faut retenir de ce chapitre~:}
\begin{itemize}
    \item \textbf{IA forensique~:} puissant pour le triage, problématique
          pour l'opposabilité. L'IA oriente~; l'humain constate et signe.
          Documenter outil, version, paramètres, échantillon de vérification.

    \item \textbf{Deepfakes~:} détection par analyse de cohérence physique,
          analyse spectrale (FFT), métadonnées et standard C2PA. La
          \loicam{2010/012} art.~72 et la \loicam{2024/017} art.~59
          couvrent les deepfakes frauduleux.

    \item \textbf{Cloud forensique~:} trois défis~: multitenancy, éphéméralité,
          multi-juridiction. La réquisition aux opérateurs cloud passe par
          des portails LEA dédiés (AWS, Azure, Google, Binance).

    \item \textbf{Mobile Afrique~:} quatre niveaux d'acquisition NIST. En
          Afrique subsaharienne, la forensique mobile est prioritaire~:
          Mobile Money, WhatsApp, SIM swapping sont les vecteurs dominants.

    \item \textbf{OSINT~:} cycle en 5~étapes, 8~outils de référence.
          La capture d'écran horodatée est obligatoire. L'infiltration en
          ligne sort du cadre OSINT et exige une autorisation judiciaire.

    \item \textbf{DICES~:} le Cameroun a un cadre législatif solide
          (\loicam{2010/012} + \loicam{2024/017}) mais manque de laboratoire
          accrédité et de jurisprudence publiée.

    \item \textbf{Roadmap 36~mois~:} Socle (Kali/Autopsy/Volatility)
          $\to$ Spécialisation (judiciaire / IR / mobile)
          $\to$ Excellence (GCFE/GCFA + contribution institutionnelle).
\end{itemize}
\end{boxresume}

\bigskip

\begin{boxexo}[title={\ding{45}\quad Exercice~5.1 --- Évaluation d'un outil IA \niveaubase}]
Un analyste utilise un outil commercial d'IA pour classifier automatiquement
500~emails dans une affaire de fraude. L'outil indique que 47~emails sont
``fortement suspects'' avec une probabilité de 0.93.

\begin{enumerate}[(a)]
    \item Quelles informations l'analyste doit-il obtenir sur l'outil
          avant de l'utiliser dans un rapport judiciaire~?
    \item Comment valide-t-il les conclusions de l'IA sur un échantillon~?
    \item Comment formule-t-il la conclusion dans son rapport pour
          qu'elle soit défendable devant un tribunal camerounais~?
    \item Formulez l'analyse CRO de cette décision d'utiliser l'IA.
\end{enumerate}
\end{boxexo}

\begin{boxexo}[title={\ding{45}\quad Exercice~5.2 --- Workflow OSINT \niveaumoyen}]
Dans une affaire de cyberharcèlement, la victime vous fournit~:
\begin{itemize}
    \item Un numéro de téléphone camerounais~: \texttt{+237 6XX XXX XXX}
    \item Un pseudo Instagram~: \texttt{@shadow\_cam237}
    \item Une photo de profil téléchargée depuis Instagram
    \item Un email reçu depuis~: \texttt{noreply@camphone-support.cm}
\end{itemize}

\begin{enumerate}[(a)]
    \item Construisez un workflow OSINT complet~: pour chaque élément,
          quel outil utilisez-vous et quelle information cherchez-vous~?
    \item Identifier les limites légales de votre collecte~: à quel
          moment devez-vous une réquisition judiciaire pour continuer~?
    \item Comment documentez-vous vos sources pour la recevabilité~?
\end{enumerate}
\end{boxexo}

\begin{boxexo}[title={\ding{45}\quad Exercice~5.3 --- Analyse DICES du Cameroun \niveauexpert}]
En vous appuyant sur le framework DICES et sur les informations du
chapitre~:
\begin{enumerate}[(a)]
    \item Pour chaque axe DICES, identifiez une force et une faiblesse
          spécifique au contexte camerounais en 2026.
    \item Proposez une action concrète et réalisable à court terme
          ($\leq$ 2~ans) pour chacune des deux lacunes les plus critiques.
    \item Dans la roadmap 36~mois, quel axe DICES est le plus influencé
          par votre choix de spécialisation~? Justifiez.
\end{enumerate}
\end{boxexo}

\bigskip

\begin{boxinfo}
\textbf{Transition vers la Partie~II.} La Partie~I vous a donné
la philosophie (Ch.~\ref{chap:phi-fond}), l'histoire
(Ch.~\ref{chap:histoire}), les affaires de référence
(Ch.~\ref{chap:gr-affaires}), les fondements mathématiques
(Ch.~\ref{chap:fond-theo}) et l'état de l'art (ce chapitre). Ces
cinq chapitres constituent le \emph{pourquoi} de la discipline. La
\textbf{Partie~II} (Chapitres~6--9) en construit le \emph{comment}~:
le Processus Investigatif Universel, les méthodologies et standards,
la chaîne de custody et les outils d'acquisition~; tout ce qu'il vous
faut pour conduire votre première investigation de A à Z.
\end{boxinfo}